% title: 试卷：Sobolev空间
% renderer: btex
% legacy-title: 用户:Solution/ 试卷: Sobolev空间
% source: https://bananaspace.org/legacy/wiki/%E7%94%A8%E6%88%B7:Solution/_%E8%AF%95%E5%8D%B7:_Sobolev%E7%A9%BA%E9%97%B4
% summary: 从香蕉空间旧站迁移的原始 bTeX 页面。
\setdisplaytitle{Sobolev空间\ 试卷}
\sethtmltitle{Sobolev空间\ 试卷}


== 2024秋 ==
(注:以下定理，例子，引理出自 
Robert A. Adams and J.F. Fournier's \textit{Sobolev Spaces}, 
Pure and Applied Mathematics Vol.140, Academic Press, 2003.)\\
以下总假定区域$\Omega$满足锥条件，$p,q\geq 1,m\geq 1,j\geq 0$.
\begin{itemize}
\item[一]证明定理4.31.
\item[二]举出四个嵌入和两个紧嵌入的反例并简单说明.(例4.41 4.42 6.12)
\item[三]解释$C^{m}(\overline\Omega),C^{m,\alpha}(\overline\Omega),L^p(\Omega),W^{m,p}(\Omega),C^\infty_0(\Omega),\mathcal{D}(\Omega),\mathcal{D}^\prime(\Omega)$,并说明它们之间的关系.
\item[四]
    (a)说明什么是锥条件.\\
    (b)证明引理4.15.\\
    (c)设$mp>n$,证明Sobolev嵌入
    \[ W^{j+m,p}(\Omega)\rightarrow C^j_B(\Omega)\]
    \[ W^{j+m,p}(\Omega)\rightarrow W^{j,q}(\Omega)\]
   其中$p\leq q \leq \infty$.
\item[五]
证明定理5.8，说明$\theta$是什么.
   
\end{itemize}